\section{Hypergeometric conversion and polynomial cases}\label{sec:hypergeom-conversion}

\subsection{Every convergent generalized hypergeometric series is a \texorpdfstring{$G$}{G}}

Put one bottom parameter $0$ in the right block and transform the others by
$a_i\mapsto1-a_i$, $b_j\mapsto1-b_j$.  In
\cref{thm:slater-right}, there is only the pole family $b_1=0$.  Its prefactor
is
\[
 \frac{\prod_i\Gamma(a_i)}{\prod_j\Gamma(b_j)},
\]
and its hypergeometric argument is $z$.  Therefore:

\begin{theorem}[Hypergeometric-to-Meijer conversion]\label{thm:hyp-to-g}
Assume no upper parameter is a nonpositive integer, no denominator parameter
is a nonpositive integer, and the series is in its convergence domain.  Common
upper/lower factors may first be cancelled.  Then
\begin{align}
{}_pF_q\!\left(\left.\begin{matrix}\mathbf a\\\mathbf b\end{matrix}\right|z\right)
&=\frac{\prod_{j=1}^{q}\Gamma(b_j)}
        {\prod_{i=1}^{p}\Gamma(a_i)}
G_{p,q+1}^{1,p}\!\left(
 \left.\begin{matrix}1-\mathbf a\\0,1-\mathbf b\end{matrix}\right|-z\right),
\label{eq:hyp-g-first}\\
&=\frac{\prod_{j=1}^{q}\Gamma(b_j)}
        {\prod_{i=1}^{p}\Gamma(a_i)}
G_{q+1,p}^{p,1}\!\left(
 \left.\begin{matrix}1,\mathbf b\\\mathbf a\end{matrix}\right|-z^{-1}\right).
\label{eq:hyp-g-second}
\end{align}
The first equality holds for every finite $z$ when $p\le q$, and for
$\abs z<1$ when $p=q+1$.  The second is the first followed by inversion and is
read on the correspondingly transported branch.
\end{theorem}
\begin{proof}
The calculation preceding the theorem turns \eqref{eq:slater-right} into
$G=(\prod\Gamma(a_i)/\prod\Gamma(b_j)){}_pF_q$.  All coefficients are equal
exactly by \cref{lem:rising}; the tails have the common modulus from
\cref{thm:hyp-conv}.  Equation \eqref{eq:hyp-g-second} follows from
\cref{thm:inversion}.  Gamma prefactors are evaluated only with pole-apartness
witnesses.
\end{proof}

\subsection{Analytic continuation and the case \texorpdfstring{$p>q+1$}{p > q + 1}}

When $p=q+1$, the series has radius one.  Formula
\eqref{eq:hyp-g-first} supplies a continuation along every certified path that
avoids the selected cut, because overlapping disks are compared by finite
coefficient equality and explicit reexpansion bounds.  The familiar principal
continuation has a cut from $1$ to $+\infty$ after the argument $-z$ and branch
conventions are translated.

For $p>q+1$, the generic hypergeometric power series diverges for every
$z\ne0$.  A $G$-contour nevertheless defines sectorial solutions of the same
formal differential equation.  This is a useful extension, but not a unique
branch-free value: different admissible contours can encode different Stokes
solutions.  The precise constructive statement is:

\begin{proposition}[Sectorial extension]
Fix a contour type, a branch sector, and all vertical or residue tail
certificates.  Then \eqref{eq:hyp-g-first}, read as a definition rather than an
equality to a convergent series, gives a primitive-recursive name solving the
formal generalized hypergeometric equation and having the prescribed
asymptotic/Frobenius data.  Two such names agree where their contours are
certifiably deformable into one another.
\end{proposition}
This qualifies the page's phrase ``the natural extension'' without weakening
the computational content.

\subsection{Terminating series}

If an upper parameter equals $-h$, the Gamma prefactor in
\eqref{eq:hyp-g-first} contains $1/\Gamma(-h)=0$, while the contour parameters
simultaneously become resonant.  The limiting product can be finite and
nonzero; multiplying a formal infinity by zero is not a proof.  The source page
repairs the problem with a sum of two $G$-functions.

Let $0\le m\le q$, $0\le n\le p$, and write
\[
 R_{m,n}(\mathbf a,\mathbf b)=
 \frac{\prod_{j=n+1}^{p}\Gamma(1-a_j)
       \prod_{j=m+1}^{q}\Gamma(b_j)}
      {\prod_{j=1}^{n}\Gamma(a_j)
       \prod_{j=1}^{m}\Gamma(1-b_j)}.
\]

\begin{theorem}[Polynomial repair formula]\label{thm:polynomial-repair}
For $h\in\N$, wherever the displayed prefactors are interpreted after the
permitted cancellations,
\begin{align}
&{}_{p+1}F_q\!\left(\left.\begin{matrix}-h,\mathbf a\\\mathbf b\end{matrix}\right|z\right)
=h!R_{m,n}(\mathbf a,\mathbf b)\,[\mathcal G_+(z)+(-1)^h\mathcal G_-(z)],
\label{eq:poly-repair}\\
&\mathcal G_+(z)=G_{p+1,q+1}^{m+1,n}\!\left(
 \left.\begin{matrix}1-\mathbf a,h+1\\0,1-\mathbf b\end{matrix}\right|
 (-1)^{p-m-n}z\right),\\
&\mathcal G_-(z)=G_{p+1,q+1}^{m,n+1}\!\left(
 \left.\begin{matrix}h+1,1-\mathbf a\\1-\mathbf b,0\end{matrix}\right|
 (-1)^{p-m-n}z\right).
\end{align}
This is a finite iPRA identity: the left side has degree $h$, and the right side
has the same coefficients through degree $h$ and zero coefficients afterward.
\end{theorem}
\begin{proof}
Expand $\mathcal G_+$ by the right residues and $\mathcal G_-$ by the left
residues, using Laurent jets if a chosen $m,n$ creates a collision.  Apply
\eqref{eq:gamma-shift-identities} and the reflection formula
$\Gamma(u)\Gamma(1-u)=\pi/\sin\pi u$.  For the coefficient of $z^r$, all
parameter-dependent constants reduce to
\[
 \frac{\poch{-h}{r}\prod_i\poch{a_i}{r}}
      {r!\prod_j\poch{b_j}{r}}.
\]
For $0\le r\le h$ this is the polynomial coefficient.  For $r>h$ the factor
$\poch{-h}{r}$ is zero; in the residue presentation this zero is the exact
cancellation of the two complementary terms.  The calculation is a finite
ring identity after multiplying by the relevant Gamma products.  Since every
coefficient after $h$ is zero, no infinite tail remains on the combined right
side.  Quantifier-free induction on $r\le h$ proves the identity in iPRA.
\end{proof}

\begin{remark}
The freedom to choose $m$ and $n$ is useful, not cosmetic.  One chooses blocks
so that no individual prefactor asks for an undefined Gamma value; if a
collision is unavoidable, the jet evaluator computes the finite combined
limit before division.
\end{remark}
